\section{Who Can Bind the State That Enforces the Rule?}
\label{sec:state-enforces}

A state can make review binding on a private developer because the state
controls something the developer needs: legal permission to operate,
access to a market, a procurement contract, or the enforcement of its
property rights. The structure becomes harder when the state is itself
the operator. The institution that writes the rule, controls the
evidence, and supplies the remedy is then the institution the rule is
supposed to constrain.

The European Union's AI Act makes the boundary visible. It can require
documentation from private providers, subject high-risk systems to
review, and impose penalties tied to global turnover
\autocite[Art.~101]{eu2024aiact}. It also excludes systems used for military,
defense, and national-security purposes from its scope. Public
authorities remain covered in other domains, but enforcement against a
member state's agencies still runs through institutions constituted by
that state. The point is not that the statute does nothing. It is that
the strongest enforcement powers face private firms, while the uses for
which the state claims its greatest discretion sit at or beyond the
instrument's edge.

That distinction matters more here than a comparison among national
regulatory models. A law governing what companies may build and sell is
not the same instrument as a law constraining what a government may do
with systems it controls directly. The first has an external enforcer.
The second asks the state to enforce a limit against itself, often in
the domains where it can classify the evidence and invoke emergency,
military, or national-security authority.

For review to bind in that setting, the protected party must be able to
initiate scrutiny without belonging to the state or its electorate. The
reviewer must be able to obtain evidence the operating agency would
prefer to withhold, issue an interim order before the disputed use is
complete, and carry the finding to a body whose remedy does not depend
on the agency's consent. An international principle, treaty obligation,
or domestic prohibition that supplies none of those powers may state
the rule correctly. It does not yet identify anyone who can make the
state obey it.

A treaty obligation of that kind exists, and it is the only one. The
Council of Europe's Framework Convention on Artificial Intelligence, the
first treaty specifically on AI, binds a party at the level of a
human-rights instrument — consistency with human rights, democracy and
the rule of law across a system's lifecycle, with the operational detail
left to domestic law \autocite{coe2024framework} — and it is not yet in
force, requiring five ratifications, at least three by Council of Europe
member states, of which the European Union deposited one in May 2026
\autocite{coe2026euratifies}. Everything else at that level is
voluntary.
